% Kompiuterijos katedros šablonas
% Template of Department of Computer Science II
% Versija 1.0 2015 m. kovas [ March, 2015]

\documentclass[a4paper,12pt,fleqn]{article}
\usepackage[unicode,colorlinks=false]{hyperref}


\usepackage[utf8x]{inputenc}
%

\usepackage[L7x]{fontenc}
\usepackage{times}
\usepackage{ucs}

 %package to switch the language
\usepackage{etoolbox}

  %set up of the page margins
\usepackage[top=2cm, bottom=2cm, left=3cm, right=1.5cm]{geometry}

 %1.1 line spacing
\linespread{1.1}


  %page numbering at the right side
\usepackage{fancyhdr}
\pagestyle{fancyplain}
\fancyhf{}
\renewcommand{\headrulewidth}{0pt} 
\fancyhfoffset[RO]{0cm}

  %to number at the bottom (exchange lines to number at the top)
\rfoot{\thepage}
  %\rhead{\thepage} %

% \usepackage[usenames,dvipsnames]{pstricks}
\urlstyle{same}
\hypersetup{
%  citecolor=Blue,
%  linkcolor=Blue,
%  urlcolor=Blue
pdfborder={0 0 0 }
}

 %for includegraphics
\usepackage{graphicx}



\usepackage[toc,page]{appendix}


\usepackage{caption}

 %for source codes
\usepackage{listings}
\lstset{commentstyle=\color{red},xleftmargin=10pt, framexleftmargin=6pt, numbersep=1mm, frame=single, numbers=left,numberstyle=\footnotesize,extendedchars=\true, inputencoding=utf8x,basicstyle=\footnotesize,extendedchars=true,
 keywordstyle=\color{black}\bfseries, breaklines=true, breakautoindent=true,framesep=8pt,linewidth=0.95\textwidth
}

 %for algorithms
\usepackage{algorithm}
\usepackage{algorithmic}
 %instead of the above two packages we can use algorithms2e
 %\usepackage[boxed,linesnumbered,vlined,slide]{algorithm2e}

 %special symbols
\usepackage{amsfonts}
\usepackage{amssymb}
\usepackage{amsmath}

 %for theorem like environments
\usepackage{amsthm}

 \usepackage{datetime}
 \renewcommand{\dateseparator}{--}


% SI system units
\usepackage{siunitx}
\sisetup{detect-all}
% Problem with fonts \SI{x.xx}{\micro\metre}, solved with updmap-sys --enable Map=utm.map
\renewcommand{\sfdefault}{uhv}
\renewcommand{\rmdefault}{utm}
\renewcommand{\ttdefault}{ucr}

% List management (itemize, etc.)
\usepackage{enumitem}

\newcommand*{\urlw}[1]{\href{#1}%
            {\nolinkurl{#1}}}
    

\numberwithin{equation}{section}


\newtoggle{inLithuanian}
 %If the report is in Lithuanian, it is set to true; otherwise, change to false
\settoggle{inLithuanian}{true}

%create file preface.tex for the preface text
%if preface is needed set to true
\newtoggle{needPreface}
\settoggle{needPreface}{false}

\newtoggle{signaturesOnTitlePage}
\settoggle{signaturesOnTitlePage}{true}


\theoremstyle{definition}
\newtheorem{definition}{\keyWordDefinition}
\newtheorem{example}{\keyWordExample}
\def\QED{\unskip\nobreak\hfill\kern5pt$\Box$}

\iftoggle{inLithuanian}{
%\usepackage[L7x]{fontenc}
\usepackage[english,lithuanian]{babel}

\newcommand{\todayiso}{\the\year \dateseparator \twodigit\month \dateseparator \twodigit\day}


\renewcommand{\today}{\number\year\space m. \space \ifcase\month\or
  sausio\or vasario\or kovo\or balandžio\or gegužės\or birželio\or
  liepos\or rugpjūčio\or rugsėjo\or spalio\or lapkričio\or
  gruodžio\fi
  \space\number\day\space d.}


 \usepackage{tocloft}
 \renewcommand\cftsecaftersnum{.} 
 \renewcommand\cftsubsecaftersnum{.} 
 \renewcommand\cftsubsubsecaftersnum{.}

 \usepackage{VUMIFKK}

 \DeclareCaptionLabelFormat{captionlt}{#2 #1}
   %smth is not fine with algorithms 
 \DeclareCaptionLabelFormat{captionltalg}{#2 #1 algoritmas}

 \usepackage{indentfirst}
 \renewcommand{\appendixtocname}{Priedai}
 \renewcommand{\appendixpagename}{Priedai}
 \renewcommand{\contentsname}{Turinys} 

 \renewcommand{\lstlistingname}{išeities kodas}
 \renewcommand{\figurename}{pav}
 \renewcommand{\tablename}{lentelė}


 \captionsetup*[lstlisting]{   
 labelsep=period,labelformat=captionlt
 }
 \captionsetup*[figure]{   
% labelsep=period,
 labelsep=space, %babel redefines pav to pav.
 labelformat=captionlt
 }
 \captionsetup*[table]{   
  labelsep=period,
  labelformat=captionlt
 }
 \renewcommand{\algorithmicrequire}{\textbf{Įvestis:}}
 \renewcommand{\algorithmicensure}{\textbf{Išvestis:}}

 \captionsetup*[algorithm]{   
 labelsep=period,labelformat=captionltalg
 }

\renewcommand{\thmhead}[3]{#2 #1#3}

}
{
%\usepackage[OT1,T1]{fontenc}
%\usepackage[L7x]{fontenc}



\usepackage[english]{babel}
\newcommand{\todayiso}{\twodigit\month \dateseparator \twodigit\day \dateseparator \the\year}
 \captionsetup*[algorithm]{   
 labelsep=period
 }
\captionsetup*[lstlisting]{   
 labelsep=period
 }
 \captionsetup*[figure]{   
 labelsep=period
 }
 \captionsetup*[table]{   
 labelsep=period
 }


}

%some kywords
 \def\keywordAbstract{\iftoggle{inLithuanian}{Santrauka}{Abstract}}
 \def\keywordAbstractOther{\iftoggle{inLithuanian}{Summary}{Santrauka}}
 \def\keyWordIntroduction{\iftoggle{inLithuanian}{Įvadas}{Introduction}}
 \def\keyWordConclusions{\iftoggle{inLithuanian}{Išvados ir rekomendacijos}{Conclusions and Recommendations}}

 \def\keyWordPreface{\iftoggle{inLithuanian}{Pratarmė}{Preface}}
 \def\keyWordAppendice{\iftoggle{inLithuanian}{Priedas}{Appendix}}
 \def\keyWordSignature{\iftoggle{inLithuanian}{parašas}{signature}}
 \def\keyWordDefinition{\iftoggle{inLithuanian}{apibrėžimas}{Definition}}
 \def\keyWordExample{\iftoggle{inLithuanian}{pavyzdys}{Example}}

\newcommand{\bothabstracts}[3]{
\setcounter{secnumdepth}{0}
\newpage
\hspace{2cm}
{\centering{\section{\keywordAbstract}}}

#1
\newpage
\hspace{2cm}
{\centering \section{\keywordAbstractOther}}

\begin{center}{\textbf{#2} }\end{center}

 #3
\setcounter{secnumdepth}{3}
}

 %non-numbered sections: #1 param: for labeling sec:#1, #2 -section title
\newcommand{\sectionWithoutNumber}[2]{\newpage
%\hspace{2cm}
\section*{#1}
\label{sec:#2}
\addcontentsline{toc}{section}{\nameref{sec:#2}}%{#3}
 }



\newcommand{\referenceSources}[1]{
\newpage
\cleardoublepage
\phantomsection
\iftoggle{inLithuanian}{
 \renewcommand{\refname}{Literatūros šaltiniai}

 \addcontentsline{toc}{section}{Literatūros šaltiniai}
 \markboth{\refname}{Literatūros šaltiniai}
 }
{

\addcontentsline{toc}{section}{References}
\markboth{References}{References}
}

\bibliographystyle{plain}
\bibliography{#1}
}



 \newcommand\authorsignature[1]{
\begin{flushright}
 \begin{minipage}[b]{0.45\textwidth}
  \centering
  \rule{\textwidth}{0.5pt}\\
   #1
  \end{minipage}
\end{flushright}
 }




 \newcommand\authorsignatures[5]{%
   \vspace{1cm}
   \authorsignature{#1}
   \ifstrequal{#2}{}{}{\vspace{0.3cm}
     \authorsignature{#2}
     \ifstrequal{#3}{}{}{\vspace{0.3cm}
      \authorsignature{#3}
      \ifstrequal{#4}{}{}{\vspace{0.3cm}
        \authorsignature{#4}
        \ifstrequal{#5}{}{}{\vspace{0.3cm}
         \authorsignature{#5}       
        }
      }
    }
} 
}

\newcommand{\authortitle}{
\iftoggle{signaturesOnTitlePage}{
\tiny{\keyWordSignature}
}{}
}

\newcommand{\depttitlepage}[8]
{
\thispagestyle{empty}
\begin{center}

\includegraphics[width=2cm]{jb_VU_zenklas}

%\vspace{-1cm}

\iftoggle{inLithuanian}
{ 
  VILNIAUS UNIVERSITETAS\\
  MATEMATIKOS IR INFORMATIKOS FAKULTETAS\\
  INFORMATIKOS INSTITUTAS\\
  KOMPIUTERINIO IR DUOMENŲ MODELIAVIMO KATEDRA
}
{
  VILNIUS UNIVERSITY \\
  FACULTY OF MATHEMATICS AND INFORMATICS \\
  INSTITUTE OF COMPUTER SCIENCE\\
  DEPARTMENT OF COMPUTATIONAL AND DATA MODELING
}

\vspace{5cm}

#1\\
\vspace{0.5cm}
\textbf{\Large #2}
\end{center}

\vspace{5cm}


\hspace{0.5\textwidth}
\begin{minipage}{0.4\textwidth}
 \begin{flushleft} 
\iftoggle{inLithuanian}
{
 \ifstrequal{#3}{}{}{Atliko:\\[5pt]}
}
{
\ifstrequal{#3}{}{}{Done by:\\[5pt]}
}

%\noindent
\begin{tabular}{@{}lr}%\setlength\tabcolsep{0pt}
\ifstrequal{#3}{}{}{#3&\hspace{2cm}\authortitle\\[5pt]}
\ifstrequal{#4}{}{}{#4&\authortitle\\[5pt]}
\ifstrequal{#5}{}{}{#5&\authortitle\\[5pt]}
\ifstrequal{#6}{}{}{#6&\authortitle\\[5pt]}
\ifstrequal{#7}{}{}{#7&\authortitle\\}
\end{tabular}

\end{flushleft}

\end{minipage}

\vspace{0.5cm}
\hspace{0.5\textwidth}
\begin{minipage}{0.4\textwidth}
 \begin{flushleft} 

\ifstrequal{#8}{}{}
{

\iftoggle{inLithuanian}
{
Vadovas:
}
{
Supervisor:
}

#8

}

\end{flushleft}

\end{minipage}


\vfill

\begin{center}
Vilnius\\
\the\year
\end{center}

\iftoggle{needPreface}{
 \sectionWithoutNumber{\keyWordPreface}{preface}
Pratarmės (Preface) informacija


\iftoggle{inLithuanian}
{
\vspace{\baselineskip}\hfill
\today
}
{
 \vspace{\baselineskip}\hfill \today
}

 \vspace{5cm}

\iftoggle{signaturesOnTitlePage}{}
{
\authorsignatures{#3}{#4}{#5}{#6}{#7}
}
}{}
\newpage
}


\begin{document}
 % #1 -report type, #2 - title, #3-7 students, #8 - supervisor
 \depttitlepage{Bakalauro darbas}{Saugus pažeidžiamumų skaitytuvas sistemų auditavimui}{Jonas Gavėnavičius } 
 {}{}{}{}% students 2-5
 {Lektorius Virgilijus Krinickij}

\tableofcontents


%keywords and notations if needed
\sectionWithoutNumber{Sutartinis terminų žodynas}{keywords}{Pateikiamas terminų sąrašas (jei reikia)}

 %both abstracts
\bothabstracts{Šiais laikais, kai pasaulis tampa vis labiau ir labiau skaitmenizuotas, iškyla opi problema – kibernetinė sauga. Vis daugiau pavyzdžių kasdieniame gyvenime matome, kai įsibraunama į sistemas, kurios yra
išnaudojamos finansiniais ar kitais tikslais. Dėl to nukenčia tiekėjai prarasdami savo reputaciją ir kapitalą, tų sistemų
vartotojai, kai jų privati informacija pavogiama ir paviešinama. Niekas nėra saugus nuo tokių situacijų. Todėl valstybės, įmonės ir korporacijos vis daugiau investuoja
į kibernetinę saugą. Kibernetinė sauga tampa vis dažnesne diskusija visuomenėje, vis daugiau
dėmesio ir resursų skiriama būtent kibernetinei saugai, stengiamasi užkirsti kelią situacijoms, kurios atneštų žalą vartotojams, įmonėms ar net šalims. Kuriami
įrankiai, kurie padeda apsisaugoti nuo tokių situacijų, analizuoja sistemas ir randa jų spragas. Įrankiai pritaiko įvairias technologijas ar metodus,
kurie perspėja apie galimą arba vykstančią ataką prieš sistemą, ištaiso rastas sistemos spragas, kol jų nerado asmenys, kurie jomis pasinaudotų.}%tex-file of abstract in original language
{Darbo pavadinimas kita kalba} %if work is in LT this title should be in English
{Nowadays, as the world becomes more and more digitized, cyber security is a major issue. As a result, states, businesses and corporations are increasingly investing in cybersecurity. Ef{}forts are being made to prevent situations that could harm consumers, businesses or even countries. Tools are developed to help prevent such situations, analyze systems and identify gaps.

The purpose of this work is to create a secure vulnerability scanning tool that can safely scan the Internet
site, its system and its files, and provide information about the scan results.

This paper introduces a secure vulnerability scanner for auditing systems and websites..
It also presents an analysis of related work that reveals best practices of other scanners and their implementation in this scanner.
In addition, the methods of vulnerability and bug analysis, their analysis and implementation in the scanner are presented.
The scanner was developed using C\# programming language, .NET Core and Angular frameworks, Docker container technology, Microsoft SQL database.
A vulnerability scanner can detect external system vulnerabilities and malicious software in web site files.
It is also capable of detecting whether a website is safe to visit, its open directories or pages that users should not have access to.
Scanner security is ensured by container technology, which ensures that malware does not enter the scanner system when scanning website files or performing other scanning operations.}%tex-file of abstract in other language


 %Introduction section: label is sec:intro
\sectionWithoutNumber{\keyWordIntroduction}{intro}
Šiame darbe kuriamas saugus įrankis, kuris neišduotų savo serverio vietos ir vartotojai pasinaudoję juo galėtų gauti išsamią informaciją apie jų internetinėje svetainėje egzistuojančias spragas
bei pažeidžiamumus ar modifikuotus failus. Darbo tikslas – sukurti saugų pažeidžiamumų skenavimo įrankį, kuris pasileistų iš tam tikros VPN technologija slepiamos vietos, skenuotų internetinę
svetainę bei jos failus, aptiktų pažeidžiamumus ar modifikuotus failus, ir pateiktų klientui informaciją apie skenavimo rezultatus. Siekiant, kad darbas būtų paklausus ir veiksmingas, keliami šie
darbo uždaviniai:
\begin{enumerate}
	\item Apžvelgti egzistuojančius panašius įrankius;
	\item Išanalizuoti dažniausiai pasitaikančius internetinių svetainių pažeidžiamumus;
	\item Išanalizuoti dažniausiai pasitaikančių pažeidžiamumų egzistuojančius atviro kodo įrankius;
	\item Išanalizuoti dažniausiai pasitaikančius internetinių svetainių skenavimo metodus.
\end{enumerate}
Kibernetinis saugumas yra svarbi kiekvienos sistemos dalis. Kadangi kiekvieną sistemą kuria
žmogus, ir į kiekvieną sistemą įeina žmogiškasis faktorius, dėl kurio sistemos beveik visada turi
pažeidžiamumų, kuriuos gali išnaudoti puolėjas siekiantis tam tikros naudos. Dėl šios priežasties
toks įrankis būtų ypač naudingas bet kokiai sistemai. Aptikus pažeidžiamumą, galima įtarti, kad
tą pažeidžiamumą gali aptikti ir nuostolio siekiantys asmenys, arba jis jau buvo aptiktas, ir tam tikri
failai buvo modifikuoti įdedant tam tikrą kodą, kuris leistų atakuojančiui asmeniui patekti į sistemą
nepastebėtam.
Dėl šių priežasčių saugus ir automatizuotas sistemų pažeidžiamumo skenavimo įrankis yra
ypač aktualus. Tačiau šio įrankio kūrimą apsunkina šios priežastys:
\begin{itemize}
	\item Įrankio kūrimas reikalauja didelio bagažo žinių;
	\item Daugumos sistemų, kurios turi pažeidžiamas vietas, išeities kodas nėra atviras, todėl kai kurios spragos negali būti aptiktos automatizuotu įrankiu. Taip pat negalima atlikti kai kurių sistemos kodo analizių, kas taip pat sumažina tokio įrankio efektyvumą;
	\item Sistemoje gali būti tiek daug skirtingų potencialių pažeidžiamų, jog jų aptikimo automatizavimas tampa problematiškas.
\end{itemize}
Šias problemas siūloma spręsti apskaičiuojant tam tikrą įrankio veiksmingumą procentaliai.
Didžiausias tokio sprendimo privalumas yra tas, kad vartotojas supras, jog tokio įrankio rezultatų
analizavimas ir pritaikymas neužtikrina visų pažeidžiamumų aptikimo ir kibernetinio saugumo
užtikrinimo.
Šio įrankio kūrimo metu, siekiant sukurti saugų pažeidžiamumų aptikimo įrankį, siekiama pagerinti sistemos saugumą, bet neužtikrinti jo. Įrankis bus skirtas tik internetinėms svetainėms ir
nebus stengiamasi jo pritaikyti ir kitoms sistemoms.



 %the main part
\newpage
\section{Susijusių darbų analizė}
\label{sec:motivation}

\subsection{Internetinių svetainių spragų skaitytuvai}
\label{sec:example}

\subsubsection{Pentest-Tools skaitytuvas}
\label{sec:data}
Pateikiamas trečio lygio poskyrio tekstas.



\newpage
\section{Sistemų auditavimas}
\label{sec:motivation}

\subsection{Elgsenos analizė}
\label{sec:example}

\subsubsection{Example}
\label{sec:data}
Pateikiamas trečio lygio poskyrio tekstas.


\subsection{Statinė spragų analizė}
\label{sec:example}

\subsubsection{Example}
\label{sec:data}
Pateikiamas trečio lygio poskyrio tekstas.

\subsection{Dinaminė spragų analizė}
\label{sec:example}

\subsubsection{Example}
\label{sec:data}
Pateikiamas trečio lygio poskyrio tekstas.

\subsection{Išorinių spragų skanavimas}
\label{sec:example}

\subsubsection{Example}
\label{sec:data}
Pateikiamas trečio lygio poskyrio tekstas.

\subsection{Vidinių spragų skanavimas}
\label{sec:example}

\subsubsection{Example}
\label{sec:data}
Pateikiamas trečio lygio poskyrio tekstas.

\subsection{Aplinkos spragų skanavimas}
\label{sec:example}

\subsubsection{Example}
\label{sec:data}
Pateikiamas trečio lygio poskyrio tekstas.



\newpage
\section{Gerosios praktikos}
\label{sec:motivation}

\subsection{Atviros kodo programinė įranga}
\label{sec:example}

\subsubsection{Privalumai}
\label{sec:data}
Pateikiamas trečio lygio poskyrio tekstas.

\subsubsection{Trūkumai}
\label{sec:data}
Pateikiamas trečio lygio poskyrio tekstas.

\subsubsection{Apibendrinimas}
\label{sec:data}
Pateikiamas trečio lygio poskyrio tekstas.

\subsection{Uždaro kodo programinė įranga}
\label{sec:example}

\subsubsection{Privalumai}
\label{sec:data}
Pateikiamas trečio lygio poskyrio tekstas.

\subsubsection{Trūkumai}
\label{sec:data}
Pateikiamas trečio lygio poskyrio tekstas.

\subsubsection{Apibendrinimas}
\label{sec:data}
Pateikiamas trečio lygio poskyrio tekstas.

\newpage
\section{Sistemų auditavimo įrankis}
\label{sec:motivation}

\subsection{Įrankio aprašas}
\label{sec:example}

Pateikiamas 4.5 poskyrio tekstas. Vienas iš šaltinių [?]. Visas turinys priklauso 4 skyriui.

\subsection{Aptiktini pažeidžiamumai}
\label{sec:example}

Pateikiamas 4.5 poskyrio tekstas. Vienas iš šaltinių [?]. Visas turinys priklauso 4 skyriui.

\subsection{Tikslumas}
\label{sec:example}

Pateikiamas 4.5 poskyrio tekstas. Vienas iš šaltinių [?]. Visas turinys priklauso 4 skyriui.

\subsection{Įgyvendinti lukeščiai}
\label{sec:example}

Pateikiamas 4.5 poskyrio tekstas. Vienas iš šaltinių [?]. Visas turinys priklauso 4 skyriui.

\subsection{Trūkumai}
\label{sec:example}

Pateikiamas 4.5 poskyrio tekstas. Vienas iš šaltinių [?]. Visas turinys priklauso 4 skyriui.


 %Conclusions section
\sectionWithoutNumber{\keyWordConclusions}{conclu}
Išvados bei rekomendacijos.


%ateities darbų gairės, planas/next steps of the work
\sectionWithoutNumber{Ateities tyrimų planas}{future}{Pristatomi ateities darbai ir/ar jų planas, gairės tolimesniems darbams....}


 %file literatureSources.bib
\referenceSources{literatureSources}



%% this part is optional
\newpage
\begin{appendices}
Dokumentą sudaro du priedai: \ref{app:a} priede  ....
\newpage
\section{Pirmojo priedo pavadinimas}
\label{app:a}
Pirmojo priedo tekstas ...

\newpage
\section{Antrojo priedo pavadinimas}
Antrojo priedo tekstas ...

\end{appendices}


\end{document}
